\chapter{Unit Roots II}
\label{cap:raiz_unitaria2}
% ── index ───────────────────────────────────────────────────────
\index{Phillips-Perron, test of|textbf}
\index{phillips_perron@\texttt{phillips\_perron()}|textbf}
\index{Zivot-Andrews, test of|textbf}
\index{zivot@\texttt{zivot()}|textbf}
\index{structural break|textbf}
\index{Newey-West, variance of|textbf}

% ─────────────────────────────────────────────────────────────────────────────
\section{Beyond the ADF: PP and Zivot-Andrews}

The previous chapter (ch.~\ref{cap:arima}) covered the ADF test, which corrects
for residual autocorrelation via lags. Two problems remain:

\begin{enumerate}
  \item The ADF is sensitive to conditional heteroscedasticity of errors (GARCH etc.)
  \item The ADF has low power in the presence of a \emph{structural break}: a
        stationary series with a level break looks like a random walk.
\end{enumerate}

\subsection*{Phillips-Perron Test (PP)}

Phillips and Perron (1988) adopt a non-parametric approach: they correct the
ADF statistic for heteroscedasticity and autocorrelation of \emph{any
form} via the Newey-West estimator:

\[
  Z_\alpha = T(\hat{\rho} - 1) - \frac{T^2 \hat{\sigma}^2}{2\hat{s}_T^2}(\hat{\lambda}^2 - \hat{\gamma}_0)
\]

where $\hat{s}_T^2$ is the long-run variance estimated with the Bartlett kernel.
The asymptotic distribution is identical to that of the ADF (Dickey-Fuller).

\subsection*{Zivot-Andrews Test (ZA)}

Zivot and Andrews (1992) relax the assumption of no structural break.
The test estimates the break date \emph{endogenously} as the value that minimizes
the $t$ statistic:

\[
  \hat{t}_B = \arg\min_\lambda \; t_\alpha(\lambda), \quad \lambda = t_B/T \in (0.15, 0.85)
\]

The tested equation includes a level shift or trend change:
\[
  \Delta y_t = c + \beta t + \gamma D U_t(\lambda) + \alpha y_{t-1} + \sum_{j=1}^k d_j \Delta y_{t-j} + \varepsilon_t
\]

$H_0$: unit root without structural break.

% ─────────────────────────────────────────────────────────────────────────────
\section{DGP Verification}

\begin{lstlisting}[caption={PP and Zivot-Andrews}]
seed(42)
// Random walk: rw_t = rw_{t-1} + e_t
// Trend-stationary series: stat_t = 0.5*t + e
// Series with break at t=150: AR(1) + level shift of 3.0

phillips_perron(df, rw)
phillips_perron(df, stat)
zivot(df, break_t)
\end{lstlisting}

\begin{lstlisting}[language={}, basicstyle=\ttfamily\small, frame=none,
  numbers=none, backgroundcolor=\color{codbg}]
PP — rw:      Zt = 77.365   CV 5%=-2.860   DOES NOT reject H₀  (unit root) ✓
PP — stat:    Zt = 2058     CV 5%=-2.860   DOES NOT reject H₀  (trend-stationary)
ZA — break_t: t  = -11.597  CV 5%=-4.800   REJECTS H₀, break at obs=147 ✓
\end{lstlisting}

The PP for \texttt{stat} (linear trend + noise) does not reject $H_0$ because the
standard test (without a deterministic trend in the auxiliary equation) cannot distinguish
stochastic from deterministic trend. The ZA correctly identifies the
structural break at $t=147$ (close to the true $t=150$) and rejects $H_0$.

% ─────────────────────────────────────────────────────────────────────────────
\section{Syntax}

\begin{lstlisting}
phillips_perron(df, var)         // PP without trend
phillips_perron(df, var, trend=true)  // PP with deterministic trend
zivot(df, var)                   // Zivot-Andrews (endogenous level break)
zivot(df, var, type="trend")     // break in trend
\end{lstlisting}

\begin{center}
\begin{tabular}{llll}
\toprule
\textbf{Test} & \textbf{Correction} & \textbf{Break?} & \textbf{CV 5\%} \\
\midrule
ADF      & Parametric lags       & No  & $-2{,}86$ \\
PP       & Non-param. Newey-West & No  & $-2{,}86$ \\
ZA       & PP + endogenous break & Yes & $-4{,}80$ \\
\bottomrule
\end{tabular}
\end{center}

\begin{nota}
  Trend-stationary series and series with structural breaks are two cases where
  the ADF/PP can lead to wrong conclusions. The ZA resolves the second; for the
  first, including a deterministic trend in the auxiliary equation
  (\texttt{trend=true}) is sufficient.
\end{nota}
